\chapter{Evaluation}\label{chap:evaluation}

This chapter demonstrates that the approach from \cref{chap:main-chapter}
works in practice.

A good evaluation chapter:
\begin{itemize}
      \item Describes clearly what was implemented and where.
      \item Explains the experimental setup with enough detail to be reproducible.
      \item Presents results honestly, including cases where the technique does not help.
      \item Interprets the numbers and connects them back to the theory.
\end{itemize}

%-----------------------------------------------------------------------
\section{Implementation}\label{sec:implementation}
%-----------------------------------------------------------------------

The implementation section should cover:
\begin{enumerate}
    \item The name of the tool extended and a brief description of its overall
          architecture (e.g., the pipeline from input parsing to proof output). 
          and any external solvers used.
    \item Which module(s) you added or changed and how they fit into the existing 
          architecture.
\end{enumerate}

%-----------------------------------------------------------------------
\section{Experimental Setup}\label{sec:setup}
%-----------------------------------------------------------------------

The experimental setup should specify:
\begin{enumerate}
    \item \textbf{Benchmark set.}
          Which set was used (e.g., TPDB, TermComp, SV-COMP, or a custom set)?
          How many examples total?  
          Are there relevant subsets by category, strategy, or problem type?

    \item \textbf{Configurations compared.}
          At minimum include the unmodified baseline and the tool with your
          extension. 
          List any competing tools with a brief description.

    \item \textbf{Hardware and time limits.}
          CPU model, amount of RAM, operating system, and per-example
          wall-clock timeout.
\end{enumerate}

%-----------------------------------------------------------------------
\section{Results}\label{sec:results}
%-----------------------------------------------------------------------

The results section should:
\begin{enumerate}
    \item Present results in a table (see \cref{tab:results}).
    \item Summarize the headline finding directly below the table, citing concrete
          numbers (additional YES answers, percentage improvement, overhead in seconds).
    \item Point out interesting subsets or individual examples where the technique
          is especially effective and where the technique is not effective.
    \item Note any surprising results to be discussed in \cref{sec:discussion}.
\end{enumerate}

\Cref{tab:results} summarizes the results of our experiments.

\begin{table}[H]
    \centering
    \caption{Experimental results on [\textit{benchmark set name}].
             \emph{YES}/\emph{NO} denotes the number of examples for which
             [\textit{the property / its negation}] was proved.
             \emph{MAYBE} means the tool returned no answer.
             \emph{TO} means the timeout was exceeded.
             \emph{Time} is the total wall-clock time in seconds.}
    \label{tab:results}
    \begin{tabular}{lccccc}
        \toprule
        \textbf{Configuration} &
        \textbf{YES} & \textbf{NO} & \textbf{MAYBE}& \textbf{Avg.~Time~(s)} \\
        \midrule
        \textit{Tool without extension} & -- & -- & -- & -- \\
        \textit{Tool with extension}    & -- & -- & -- & -- \\
        \bottomrule
    \end{tabular}
\end{table}

%-----------------------------------------------------------------------
\section{Discussion}\label{sec:discussion}
%-----------------------------------------------------------------------

The discussion should address the following questions:
\begin{enumerate}
    \item Do the results match theoretical expectations from \cref{chap:main-chapter}?
          If not, explain why.
    \item On which types of examples does the technique help most? 
          Identify structural properties of the benchmark examples that explain the gain.
    \item Are there cases where the technique does not help or is slower?
          Explain the reason and relate it to the theoretical properties.
    \item What are the main limitations and how could they be overcome?
          (This can bridge to the future work discussion in \cref{chap:conclusion}.)
\end{enumerate}
